\documentclass[11pt, a4paper]{report}

\usepackage[utf8]{inputenc}
\usepackage[T1]{fontenc}
\usepackage[french]{babel}

% Appel des packages dans le dossier preambule
\def\packagepath{../../../preambule}   % path du package princial

\usepackage{\packagepath/page_layout} % <--- Nouveau
\usepackage{\packagepath/config_info}
\usepackage{\packagepath/cadres_cours}
\usepackage{\packagepath/zones_reponse}
\usepackage{\packagepath/config_ds}
\usepackage{\packagepath/config_maths}

% --- PILOTAGE ---
%\def\visibleornot{visible}    % visible or invisible
\ifdefined\visibleornot\else\def\visibleornot{invisible}\fi


\def\documentpath{./Sources_Latex}
\graphicspath{{./Images/}}


\def\ptsexoA{2}
\def\ptsexoB{2}
\def\ptsexoC{2}
\def\ptsexoD{2}
\def\ptsexoE{2}
\def\ptsexoF{2}

\begin{document}

\def\level{Premiere}  % Classe à droite
\def\course{NSI}
\def\chaptername{Représentation des nombres flottants et des chaines de caractères} % Titre au centre
\def\docname{DS\quad}        % Identifiant à gauche

\pagestyle{DS_LP}

\NomPrenomNote{}

\bigskip




%%=============================================================
\exods{\ptsexoA}{\quad}

\begin{enumerate}
    \item Convertir le nombre entier naturel \texttt{52} en binaire
    
    
    \begin{blocvisible}[height=3.]
            {\quad }
            
            Le nombre 52 peut être décomposé en puissances de 2 de la manière suivante :
\[
52 = 32 + 16 + 4 = 2^5 + 2^4 + 2^2 = 101100_2
\]

    \end{blocvisible}
    
    \medskip

    \item Convertir le nombre binaire \texttt{101101} en entier naturel (positif)

    \begin{blocvisible}[height=3.]
            {\quad }
            
            Le nombre binaire 101101 peut être converti en entier naturel en utilisant les puissances de 2 :
\[
101101 = 2^5 + 2^4 + 2^0 = 32 + 16 + 1 = 53
\]

    \end{blocvisible}

\end{enumerate}


%%=============================================================
\exods{\ptsexoB}{Convertir le nombre $5,75$ en binaire. Justifier votre calcul.}

\begin{blocvisible}[height=5.5]
            {\quad }
            
            Le nombre 5,75 peut être décomposé en partie entière et partie décimale :
$5,75 = 5 + 0,75$

La partie entière 5 peut être convertie en binaire :
$5 = 4 + 1 = 2^2 + 2^0 = 101_2$


La partie décimale 0,75 peut être convertie en binaire en multipliant successivement par 2 :
$0,75 \times 2 = 1,5 \quad \Rightarrow \quad 1$


$0,5 \times 2 = 1 \quad \Rightarrow \quad 0,5 \times 2 = 1,0 \quad \Rightarrow \quad 0$

Ainsi, la partie décimale 0,75 correspond à 0,11 en binaire. En combinant les deux parties, on obtient :
$
5,75 = 101,11_2
$
    
\end{blocvisible}

\exods{\ptsexoC}{Convertir le nombre binaire $1101,101$ en décimal. Justifier votre calcul.}

\begin{blocvisible}[height=6.5]
            {\quad }
            
            Le nombre binaire 1101,101 peut être décomposé en partie entière et partie décimale :
$1101,101 = 1101 + 0,101$

La partie entière 1101 peut être convertie en décimal en utilisant les puissances de 2 :        

$1101 = 2^3 + 2^2 + 2^0 = 8 + 4 + 1 = 13$

La partie décimale 0,101 peut être convertie en décimal en utilisant les puissances de 2 négatives :


$0,101 = 1 \times 2^{-1} + 0 \times 2^{-2} + 1 \times 2^{-3} = 0,5 + 0 + 0,125 = 0,625$ 

Ainsi, en combinant les deux parties, on obtient :

$1101,101 = 13 + 0,625 = 13,625$
    
\end{blocvisible}


\exods{\ptsexoD}{On suppose que l'on dispose des deux fonctions suivantes:}
\begin{itemize}
    \item \verb|conv_partie_entiere(n)| : qui prend en paramètre un entier naturel $n$ et renvoie sa représentation binaire sous forme de chaîne de caractères.
    \item \verb|conv_partie_decimale(x)| : qui prend en paramètre un nombre décimal $x$ compris entre 0 et 1, et renvoie sa représentation binaire sous forme de chaîne de caractères.    
\end{itemize}

\begin{CodePythontexAlt}[TaillePolice=\large, Largeur=1\linewidth,Centre,Lignes=true,PremLigne=1]{}
assert conv_partie_entiere(5) == "101"
assert conv_partie_entiere(13) == "1101"
assert conv_partie_decimale(0.9375) == "1111"
assert conv_partie_decimale(0.8125) == "1110"
\end{CodePythontexAlt}

Compléter la fonction suivante qui permet de convertir un nombre décimal en binaire en utilisant les fonctions précédentes.

\begin{CodePythontexAlt}[TaillePolice=\large, Largeur=1\linewidth,Centre,Lignes=true,PremLigne=1]{}
def conv_decimal_binaire(nombre):
    ''' Convertit un nombre decimal en binaire.
    nombre : un nombre decimal (float)
    '''
    










    

assert conv_decimal_binaire(5.75) == "101,11"
assert conv_decimal_binaire(13.625) == "1101,101"

\end{CodePythontexAlt}


\exods{\ptsexoE}{\quad}

On suppose que l'on dispose d'un dictionnaire \verb|morse| suivant:

\begin{CodePythontexAlt}[TaillePolice=\large, Largeur=1\linewidth,Centre,Lignes=true,PremLigne=1]{}
morse = {'A': '.-', 'B': '-...', 'C': '-.-.', 'D': '-..', 'E': '.', 'F': '..-.',
         'G': '--.', 'H': '....', 'I': '..', 'J': '.---', 'K': '-.-', 
         'L': '.-..', 'M': '--', 'N': '-.', 'O': '---', 'P': '.--.', 
         'Q': '--.-', 'R': '.-.', 'S': '...', 'T': '-', 'U': '..-',
         'V': '...-', 'W': '.--', 'X': '-..-', 'Y': '-.--', 'Z': '--..'}

\end{CodePythontexAlt}


\begin{enumerate}
    \item Ecrire une fonction \verb|conv_morse(texte)| qui prend en paramètre une chaîne de caractères \verb|texte| et qui renvoie la traduction de ce texte en code Morse. On supposera que le texte ne contient que des lettres majuscules et des espaces.
    
    On respectera les conventions suivantes:
    \begin{itemize}
        \item Chaque lettre est séparée par un espace en Morse.
        \item Les espaces entre les mots sont représentés par " / " en Morse.
    \end{itemize}
    \begin{CodePythontexAlt}[TaillePolice=\large, Largeur=1\linewidth,Centre,Lignes=true,PremLigne=1]{}
def conv_morse(texte):
    ''' Convertit un texte en code Morse.
    texte : une chaîne de caractères contenant uniquement des lettres 
    majuscules et des espaces
    Les espaces sont représentés par " / " en Morse.
    '''
    
    

    








    # Tests
assert conv_morse("PYTHON") == ".--. -.-- - .... --- -."
assert conv_morse("HELLO WRLD") == ".... . .-.. .-.. --- / .-- .-. .-.. -.."
    
    \end{CodePythontexAlt}
\end{enumerate}
\end{document}